\chapter{Conclusão}
\label{cap:conclusao}

A digitalização direcionada e arquitetada especificamente para o acompanhamento odontológico contínuo mostrou-se um caminho promissor. A proposta não se resume a evitar a perda física de papéis: trata-se de introduzir uma cultura analítica, orientada a dados, no ambiente de ensino e prática clínica da UEPG, sem transferir à máquina a responsabilidade pelo diagnóstico.

\todo{retomar cada objetivo específico e indicar como foi atingido, com os principais números dos experimentos.}

\section{Trabalhos futuros}

\begin{itemize}
	\item Implantação em atendimentos reais após aprovação ética e adoção de termo de consentimento para gravação de áudio;
	\item Assinatura digital com certificação ICP-Brasil, permitindo a eliminação do prontuário em papel;
	\item Migração da ficha para o DC/TMD, aproveitando o esquema versionado;
	\item Preenchimento do questionário de autorrelato pelo próprio paciente em \textit{tablet}, na sala de espera;
	\item Interoperabilidade com a Rede Nacional de Dados em Saúde por meio do padrão HL7 FHIR;
	\item Ampliação para outras ligas e clínicas do curso de Odontologia da UEPG.
\end{itemize}
